\documentclass{article}
\usepackage[T1]{fontenc}
\usepackage{lmodern}
\usepackage{graphicx}
\usepackage{amsmath,amssymb}
\usepackage[a4paper,margin=2cm]{geometry}
\usepackage[absolute,overlay]{textpos}
\setlength{\TPHorizModule}{1mm}
\setlength{\TPVertModule}{1mm}
\usepackage[indonesian]{babel}
\usepackage{hyperref}
\setlength{\emergencystretch}{2em}
% unit: Halaman 19

\begin{document}

% === Halaman 19 (python/native) ===

Pertukaran kunci Diffie-Hellman untuk tiga 

pihak (Alice, Bob, dan Carol)

1. Alice, Bob, dan Carol menyepakati p dan g.

2. Alice, Bob, dan Carol membangkitkan kunci privat masing-masing, a, b, dan c.

3. Alice menghitung ga mod p dan mengirimkannya kepada Bob.

4. Bob menghitung (ga)b mod p = gab mod p dan mengirimkannya kepada Carol.

5. Carol menghitung  K = (gab)c mod p = gabc mod p.

6. Bob menghitung gb mod p dan mengirimkannya kepada Carol.

7. Carol menghitung (gb)c mod p = gbc mod p dan mengirimkannya kepada Alice.

8. Alice menghitung K = (gbc)a mod p = gbca mod p = gabc mod p.

9. Carol menghitung gc mod p dan mengirimkannta kepada Alice.

10. Alice menghitung  (gc)a mod p = gca mod p dan mengirimkannya kepada Bob.

11. Bob menghitung K = (gca)b mod p = gcab mod p = gabc mod p.

12. Sekarang Alice, Bob, dan Carol sudah memiliki kunci rahasia yang sama, yaitu K

Rinaldi M/II4020 Kriptografi

19

\end{document}
